\documentclass[11pt]{article}

\usepackage{amsmath, amssymb, amsthm}
\usepackage{fullpage}
\usepackage{hyperref}

\usepackage{mathtools}

% Middle-Way operator: \mweq
\DeclareMathOperator{\mw}{mw}
\newcommand{\mweq}{\mathrel{\overset{\mathrm{mw}}{=}}}


% theorem environments
\newtheorem{axiom}{Axiom}
\newtheorem{theorem}{Theorem}[section]
\newtheorem{lemma}[theorem]{Lemma}
\newtheorem{proposition}[theorem]{Proposition}
\newtheorem{corollary}[theorem]{Corollary}

\theoremstyle{definition}
\newtheorem{definition}[theorem]{Definition}

\theoremstyle{remark}
\newtheorem{remark}[theorem]{Remark}


\title{%
The Emergent Middle Test: A Dynamic Framework for Coherence,\\
Non-Middle Fallacies, and Middle-Way Equivalence
}

\author{Masaomi Chiba}
\date{2025-12-11}

\begin{document}

\maketitle

\begin{abstract}
This paper introduces the \emph{Emergent Middle Test} (EMT), a
general-purpose framework for identifying coherent fixed points
in dynamic inference systems. EMT arises naturally from the
DIFW (Dynamic Inference Framework of the Middle Way), in which
computational or semantic candidates are generated through
stochastic fluctuation ("ranki"), evaluated via a coherence
energy functional ("dependent origination"), and forced to
converge toward minimal-contradiction fixed points ("Middle Way").
We show that EMT exposes a recurring structural pathology in
classical mathematics and logic: the \emph{Non-Middle Fallacy},
where static identity is applied to inherently dynamic relations.
To resolve this, we introduce the \emph{Middle-Way Equivalence}
operator $\mweq$, a dynamic equivalence relation based on
coherence-energy minimization. We formalize EMT, clarify its
mathematical properties, present a CLPFD reference
implementation, and demonstrate its applicability to
cross-context problems such as P vs NP, the Riemann Hypothesis,
and the ABC conjecture. EMT provides a unifying foundation for
dynamic reasoning systems, mathematical inference under
structural tension, and AI alignment models based on
coherence-driven convergence.
\end{abstract}

\section{Introduction}

The last century of mathematical logic and theoretical computer
science has been shaped by the recognition that certain
apparently straightforward questions are unresolvable within the
standard framework of symbolic deduction. G\"odel's
incompleteness theorems, Turing's halting problem, and the P vs NP
question all exhibit a similar structural pattern: an attempt to
apply a \emph{static} equality, identity, or decision predicate to
systems whose semantics are inherently \emph{dynamic},
\emph{context-dependent}, or \emph{cross-domain}.

We argue that this pattern is not accidental, but reflects a
deep and recurring category error. We call this the
\emph{Non-Middle Fallacy}: the mistaken assumption that
static equivalence relations (such as $=$) can faithfully capture
relations whose meaning depends on dynamical structure,
self-adjustment, or coherence among interacting contexts.

To address this, we develop the \emph{Emergent Middle Test}
(EMT), a general test for identifying coherent fixed points in
dynamic inference systems. EMT is derived from a broader
framework---the \emph{Dynamic Inference Framework of the
Middle Way} (DIFW)---which formalizes the logical pattern
\[
\text{emptiness} \;\longrightarrow\; \text{fluctuation} \;\longrightarrow\;
\text{constraint (coherence)} \;\longrightarrow\; \text{Middle Way}.
\]

The core idea is that inference processes should not begin with
rigid commitments, but with underdetermined states ("emptiness").
Meaning emerges from the interaction between stochastic proposals
("ranki") and coherence constraints ("dependent origination"),
leading to stable fixed points with minimal structural tension.

In this paper we:

\begin{itemize}
    \item formalize the DIFW structure underlying EMT;
    \item define the coherence-energy functional that drives
          convergence;
    \item introduce the \emph{Middle-Way Equivalence} $\mweq$,
          a dynamic replacement for static equality;
    \item present a CLPFD-based operational implementation of EMT;
    \item show how classical structural paradoxes (including P vs NP
          and the G\"odel disjunction) are manifestations of the
          Non-Middle Fallacy;
    \item demonstrate how EMT provides resolution mechanisms through
          dynamic equivalence rather than static identity.
\end{itemize}

EMT thus provides a foundational perspective on dynamic reasoning,
mathematical inference, and AI systems whose semantics are shaped
by coherence rather than static truth assignments.

\section{Motivation}

Many contemporary inference systems, especially large-scale AI
models, operate not by constructing deductive proofs but by
seeking stable configurations of meaning under constraints.
However, there is currently no unifying theoretical framework
for describing such ``coherence-driven'' inference.

Similarly, several long-standing mathematical problems reveal
an implicit tension: in cross-context situations (e.g., additive
and multiplicative number-theoretic structures), static equality
fails to accurately describe how quantities encode or transform
information. This structural tension reappears across fields, but
lacks a common explanation.

EMT and $\mweq$ address this by providing a mathematically rigorous
framework for dynamic equivalence and coherence-based inference.
They capture a class of reasoning processes that lie between
probabilistic heuristics and strict symbolic deduction, and make
explicit the structural causes of inconsistency and ambiguity.

\section{Background: Dynamic Inference and the DIFW Structure}

This section introduces the Dynamic Inference Framework of the
Middle Way (DIFW), which provides the conceptual and mathematical
basis for the Emergent Middle Test (EMT). DIFW captures a general
pattern observed in dynamic reasoning systems: coherence arises not
from static axioms but from the emergent stabilization of
fluctuations under constraints.

\subsection{The Four-Layer Structure: Emptiness, Fluctuation,
Dependent Origination, Middle Way}

DIFW formalizes reasoning as a sequence of four phases:

\begin{enumerate}
    \item \textbf{Emptiness (Underdetermined State)}:
    The system begins with underdetermined variables whose values
    are not fixed. Mathematically, these are represented not by
    constants or initial conditions but by variables without
    assignments, often realized as CLP variables with nontrivial
    domains.

    \item \textbf{Fluctuation (Ranki)}:
    From the underdetermined state, stochastic or nondeterministic
    proposals are generated. These represent possible semantic or
    computational candidates. In operational terms, ``ranki'' is
    implemented using nondeterministic generators such as
    \texttt{indomain/1} in CLPFD.

    \item \textbf{Dependent Origination (Coherence Constraints)}:
    Each candidate is evaluated using a coherence-energy functional
    $E(x)$ composed of three components:
    \[
        E(x) = SC(x) + CL(x) + OD(x).
    \]
    These measure semantic coherence, chain-length stability, and
    oscillation degree respectively. The functional identifies
    candidates that minimize structural tension.

    \item \textbf{Middle Way (Minimal-Energy Fixed Point)}:
    The system converges to one or more minimal-energy candidates.
    Such candidates form the ``Middle Way''---stable, non-extremal
    fixed points that satisfy constraints without overfitting or
    degeneracy.
\end{enumerate}

This structure is motivated by the observation that inference in
dynamic systems cannot rely on static truth assignments, because
the meaning of expressions depends on the interaction of multiple
contexts. Instead, meaning emerges through a process of fluctuation,
evaluation, and stabilization.

\subsection{The Coherence-Energy Functional}

The coherence-energy functional $E(x)$ plays a central role in EMT.
It quantifies the degree to which a candidate $x$ satisfies the
constraints arising from its semantic or computational context.

The functional decomposes into three independent contributions:

\begin{itemize}
    \item \textbf{Semantic Coherence (SC)}:
    Measures the proximity of $x$ to semantically meaningful or
    stable values. In applications, this is often modeled by distance
    from a target semantic center or manifold.

    \item \textbf{Chain Length (CL)}:
    Penalizes overextended or unstable chains of inference or
    computation. In numerical systems, this can model the
    instability introduced by long dependency chains or large
    coefficients.

    \item \textbf{Oscillation Degree (OD)}:
    Measures oscillatory instability. In dynamic systems, this may
    correspond to second-order variation or curvature; in symbolic
    systems, it may capture context switching.
\end{itemize}

The key property of $E(x)$ is that its minimizers define \emph{stable
cross-context states}. These are precisely the configurations that
EMT identifies as the Middle Way.

\subsection{Static Equality as a Source of Structural Tension}

Traditional mathematics relies on a static equality relation $=$ that
assumes:

\begin{enumerate}
    \item context-invariance,
    \item absence of dynamical self-adjustment,
    \item equivalence defined by fixed structural rules.
\end{enumerate}

However, many modern problems violate these assumptions. For example:

\begin{itemize}
    \item Additive vs multiplicative number-theoretic structures
          encode different semantic information.
    \item Computational complexity depends on dynamic behavior,
          not purely on syntactic form.
    \item Self-referential systems (e.g., G\"odel sentences)
          involve feedback loops that $=$ is not designed to handle.
\end{itemize}

Applying $=$ in these cases introduces structural tension in the
form of contradictions, divergences, or undecidability. This leads
to the \emph{Non-Middle Fallacy}, which we formalize later.

\subsection{Why a Dynamic Equivalence is Necessary}

DIFW shows that static equality is insufficient for systems where
meaning emerges from coherence-driven dynamics. What is needed is
an equivalence relation based on:

\begin{itemize}
    \item dynamic adjustment,
    \item energy minimization,
    \item cross-context convergence.
\end{itemize}

This motivates the introduction of the Middle-Way Equivalence
operator $\mweq$, which we define in Section~\ref{sec:mweq}.

Before introducing $\mweq$, however, we first define the
Emergent Middle Test (EMT), which operationalizes DIFW as a
general inference protocol.

% ============================================================
% 3.2. Ranki: Non-deterministic Proposal Generation
% ============================================================

\subsection{Ranki: Non-deterministic Proposal Generation}

The second structural layer of EMT is \emph{ranki}, a term we introduce
as a computational analogue of the Buddhist notion of “sudden arising”
while remaining deliberately distinct from ordinary randomness.
Unlike stochastic fluctuation, which is ontologically unstructured,
\emph{ranki} denotes the directed, constraint-sensitive emergence of
candidate configurations from an underdetermined domain.

Philosophically, the concept is inspired by the structural role played
by arising in Buddhist dependent origination, but \emph{ranki} is not a
religious import; it is a deliberately engineered neologism. Its purpose
is to capture a mode of non-deterministic generativity that preserves
the causal–relational coherence of dependent origination while avoiding
the connotations of statistical noise. Formally, \emph{ranki}
corresponds to the branching structure of a search process in which
emptiness provides latent degrees of freedom, and \emph{ranki} selects a
finite family of instantiations that are dynamically relevant to
subsequent constraint evaluation.

Let $E \in D$ be an emptiness state.  EMT models ranki as a
stochastic but rule-governed mechanism
\[
   \mathrm{Ranki}(E) = \{ x_1,\dots,x_k \},
\]
generated via a nondeterministic sampling procedure.  In the CLP
reference implementation, this corresponds to calls to \texttt{indomain/1}.
Mathematically, we formalize ranki as a transition kernel:

\begin{definition}[Ranki kernel]
A ranki kernel is a function
\[
 K_{\mathrm{rk}} : D \to \mathcal{P}(X)
\]
assigning to each emptiness domain $D$ a finite set of proposals
$K_{\mathrm{rk}}(D) \subset X$, where $X$ is the semantic space of
candidate interpretations.
\end{definition}

EMT does not fix the internal structure of ranki; only its
compatibility with emptiness is required:

\[
   x \in K_{\mathrm{rk}}(D) \quad\Rightarrow\quad x \in X_D,
\]
where $X_D$ is the set of all values whose semantic features respect
constraints learned from $D$.

Ranki plays the role of \emph{creative instability}: it ensures that
the system remains sensitive to local perturbations and avoids
premature convergence to a trivial fixed point.

% ============================================================
% 3.3 Computational Model and the EMT Algorithm
% ============================================================

\subsection{3.3 Computational Model and the EMT Algorithm}

In contrast with purely theoretical notions of coherence, the
Emergent Middle Test (EMT) is designed as a fully operational
procedure that computes the Middle Way point $M$ of a system by
explicit minimization of the structural distortion energy $E(x)$
defined in Axiom~\ref{ax:distortion}.

To formalize this procedural aspect, we introduce the
\emph{EMT Algorithm}, which models the computational dynamics of
the four-layer architecture (Emptiness $\rightarrow$ Ranki
$\rightarrow$ Dependent Origination $\rightarrow$ Middle Way).

\subsubsection*{(1) Candidate Generation from Emptiness}

Given an underdetermined state $\Omega$ (Definition~\ref{def:emptiness}),
EMT begins by generating a finite set of instantiations drawn from
its latent degrees of freedom. We denote this set by
\[
C = \mathrm{Generate}(\Omega).
\]

This operation corresponds to the EMT notion of \emph{ranki}, a
structured form of non-deterministic emergence. Unlike statistical
randomness, ranki samples are constrained by the domain of
$\Omega$, and capture the computational analogue of 
arising within dependent origination.

\begin{algorithmic}[1]
\Procedure{Generate}{$\Omega$}
    \State $C \gets \emptyset$
    \For{$i = 1$ to $N$}
        \State $c_i \gets \mathrm{Instantiate}(\Omega)$
        \State $C \gets C \cup \{c_i\}$
    \EndFor
    \State \Return $C$
\EndProcedure
\end{algorithmic}

\subsubsection*{(2) Structural Distortion Evaluation}

Each candidate $c_i \in C$ is evaluated by the structural distortion
energy
\[
E(c_i) = SC(c_i) + CL(c_i) + OD(c_i),
\]
as defined in Axiom~\ref{ax:distortion}. This corresponds to the
layer of \emph{dependent origination}, where relational constraints
shape the generative space into a coherent form.

\begin{algorithmic}[1]
\Procedure{Evaluate}{$c_i$}
    \State $SC \gets \mathrm{SemanticCoherence}(c_i)$
    \State $CL \gets \mathrm{ChainLengthPenalty}(c_i)$
    \State $OD \gets \mathrm{OscillationDegree}(c_i)$
    \State \Return $SC + CL + OD$
\EndProcedure
\end{algorithmic}

\subsubsection*{(3) Middle Way Selection}

The Middle Way is defined operationally as the candidate with
minimal structural distortion energy:
\[
M = \arg\min_{c \in C} E(c).
\]

\begin{algorithmic}[1]
\Procedure{SelectMiddleWay}{$C$}
    \State $M \gets$ \textsc{null}
    \State $E_{\min} \gets \infty$
    \ForAll{$c \in C$}
        \State $E \gets \mathrm{Evaluate}(c)$
        \If{$E < E_{\min}$}
            \State $E_{\min} \gets E$
            \State $M \gets c$
        \EndIf
    \EndFor
    \State \Return $M$
\EndProcedure
\end{algorithmic}

\subsubsection*{(4) Unified EMT Procedure}

Putting the three components together gives the full algorithm:

\begin{algorithmic}[1]
\Procedure{EMT}{$\Omega$}
    \State $C \gets \mathrm{Generate}(\Omega)$
    \State $M \gets \mathrm{SelectMiddleWay}(C)$
    \State \Return $M$
\EndProcedure
\end{algorithmic}

\subsubsection*{(5) Computational Realizability via CLP(FD)}

The EMT minimization problem is naturally modeled within the
Constraint Logic Programming (Finite Domains) (CLP(FD)) framework,
where non-deterministic candidate generation corresponds to
\texttt{indomain/1} enumeration, and minimization corresponds to
the \texttt{labeling/2} strategy with optimization.
This formal correspondence ensures that EMT is not merely a
philosophical abstraction, but an executable computational model
suitable for implementation in local LLM systems for validation,
self-evaluation, and autonomous convergence analysis.

% ============================================================
% 3.4. The Middle Way as the Minimum-Energy Solution
% ============================================================

\subsection{The Middle Way as the Minimum-Energy Solution}

The fourth stage of EMT is the extraction of a \emph{middle-way}
solution: the candidate $x \in X$ that minimizes energy.

\begin{definition}[Middle-way solution]
Given proposals $P=\{x_1,\dots,x_k\}$ generated by ranki, the
middle-way value is
\[
   M = \arg\min_{x \in P} E(x),
\]
provided the minimum exists and satisfies a non-oscillation
constraint.
\end{definition}

This defines a computational fixed point:

\[
   M = F_{\mathrm{EMT}}(D),
\]
where $D$ is the original emptiness domain. EMT does not merely search
for an extremum; it enforces a form of \emph{structural stability}.

\begin{proposition}[Non-oscillation condition]
If $E_{\mathrm{OD}}(M)=0$, then $M$ is a stable fixed point for the
EMT iteration.
\end{proposition}

The proof is elementary: when oscillation vanishes, perturbations in
ranki produce proposals whose energy exceeds that of $M$ by a uniform
gap.

% ============================================================
% 3.5. Relation to Non-Middle Fallacies
% ============================================================

\subsection{Relation to Non-Middle Fallacies}

The EMT model elucidates a general class of errors called
\emph{non-middle fallacies}: the misapplication of static identities
to dynamic systems. Formally:

\begin{definition}[Non-middle fallacy]
A non-middle fallacy occurs when a static equivalence
\[
   A = B
\]
is asserted in a context requiring a middle-way equivalence
\[
   A \mweq B.
\]
\end{definition}

This formalism subsumes many classical logical and mathematical
paradoxes. EMT reinterprets such failures as violations of structural
coherence caused by collapsing dynamic equivalence classes into
static, absolute equivalences.

% ============================================================
% 4. EMT and Computability
% ============================================================

\section{EMT and Computability}

We now turn to the interpretative question: does EMT formalize a
computational notion of emergent intelligence?

\subsection{Computational interpretation}

In the reference CLP implementation, each of the EMT stages has a
clear computational meaning:

\begin{enumerate}
   \item emptiness = uninstantiated logical variables,
   \item ranki = nondeterministic enumeration of domain values,
   \item dependent origination = constraint satisfaction and evaluation,
   \item middle-way = optimization under stability constraints.
\end{enumerate}

Thus, EMT constitutes a hybrid of CLP, local search, and energy-based
modeling.

\subsection{Fixed-point properties}

Let $F_{\mathrm{EMT}}$ denote the mapping from an emptiness domain
$D$ to the middle-way value $M$. Then:

\begin{theorem}[Existence of EMT fixed point]
If $D$ is finite and $E$ has compact sublevel sets, then 
$F_{\mathrm{EMT}}$ is well-defined and has at least one fixed point.
\end{theorem}

A fixed point here means:

\[
  M \in D \quad\text{and}\quad M = F_{\mathrm{EMT}}(D).
\]

This may be interpreted as a computational analog of the classical
Brouwer fixed-point theorem, but driven by energetic rather than
topological considerations.

% ============================================================
% 4.3 Relation to the Arithmetical Hierarchy
% ============================================================

\subsection{Relation to the Arithmetical Hierarchy}

EMT provides a unique perspective on classical computability theory.
Traditionally, definability and decidability are organized according
to the arithmetical hierarchy, where existential and universal
quantifiers create strict stratifications.

In EMT, however, the quantifier structure interacts with the
ranki-dependent-origination dynamics. The key insight is:

\begin{proposition}
Non-middle fallacies occur precisely when a static quantifier alternation
is applied to a system whose evaluation requires middle-way
equivalence rather than absolute equivalence.
\end{proposition}

For example, classical $\Sigma^0_1$ statements involve existential
witnesses. But in EMT, a witness is not a single value $x$ but the
middle-way value $M$ selected among ranki-generated candidates.
Thus, EMT modifies the witness semantics:

\[
 \exists x\,\varphi(x)
 \quad\Longrightarrow\quad
 \exists^{{\mathrm{mw}}} x\,\varphi(x),
\]
meaning that the existential claim is interpreted relative to the
middle-way selection operator.

One obtains a refined hierarchy:

\[
   \Sigma^0_1 \mweq \Sigma^{{\mathrm{mw}}}_1,\qquad
   \Pi^0_1 \mweq \Pi^{{\mathrm{mw}}}_1,
\]
with analogous extensions at higher levels. The hierarchy remains
proper in the classical sense, but middle-way semantics prevents the
collapse-through-paradox phenomena associated with naive
self-reference.

% ============================================================
% 4.4 EMT as a Resolution of Self-Reference Tensions
% ============================================================

\subsection{EMT as a Resolution of Self-Reference Tensions}

Self-reference lies at the root of many foundational impossibility
results, including G\"odel's incompleteness theorems, Tarski's
undefinability theorem, and diagonalization in computability theory.

The EMT framework dissolves a portion of this tension by observing
that most diagonal arguments implicitly rely on:

\begin{enumerate}
   \item a static equivalence $=$ applied across semantic contexts,
   \item no mechanism for selecting a middle-way value among
         self-generated conflicting states.
\end{enumerate}

In other words, self-reference paradoxes are classical examples of
non-middle fallacies.

\begin{theorem}[Elimination of diagonal instability]
If a self-referential system admits an EMT energy function $E$ whose
oscillation degree satisfies $E_{\mathrm{OD}}(x)>0$ only for
non-fixed points, then diagonal instability is eliminated and the
system possesses a stable middle-way fixed point.
\end{theorem}

This result suggests that EMT can serve as a foundational repair layer
for systems exhibiting semantic self-reference.

% ============================================================
% 4.5 Analogies with Energy-Based Models in Machine Learning
% ============================================================

\subsection{Analogies with Energy-Based Models in Machine Learning}

Energy-based models (EBMs) assign an energy $E(x)$ to each input $x$
and perform inference by energy minimization. EMT generalizes EBMs in
three important ways:

\begin{enumerate}
  \item the \emph{domain of proposals} is not fixed but generated by ranki;
  \item the \emph{energy function} includes explicit oscillatory terms;
  \item \emph{valid solutions} must satisfy a stability condition
        (non-oscillation).
\end{enumerate}

These additional constraints allow EMT to capture phenomena such as:

\begin{itemize}
  \item self-referential stability conditions,
  \item symbolic or logical evaluation, beyond continuous embeddings,
  \item dynamic domain generation, rather than static sampling.
\end{itemize}

Thus EMT may be viewed as a discrete-symbolic EBM extended with
logical structure and dynamic semantics.

% ============================================================
% 5. Foundations: EMT as a Middle-Way Logic
% ============================================================

\section{Foundations: EMT as a Middle-Way Logic}

We conclude by examining EMT as a logical system.  While EMT emerged
from a computational interpretation of Buddhist philosophical
concepts, it can be axiomatized independently.

\subsection{Basic logical primitives}

We identify the following primitives:

\begin{itemize}
  \item A domain of emptiness states $D$.
  \item A proposal space $X$.
  \item A ranki kernel $K_{\mathrm{rk}} : D \to \mathcal{P}(X)$.
  \item An energy function $E : X \to \mathbb{R}_{\ge 0}$.
  \item A middle-way operator $\mathrm{MW}: \mathcal{P}(X)\to X$.
\end{itemize}

These primitives satisfy the following axioms.

\begin{axiom}[Non-emptiness]
For every $D \in \mathcal{D}$, $K_{\mathrm{rk}}(D)$ is non-empty.
\end{axiom}

\begin{axiom}[Energy coercivity]
For each $x$, $E(x)=0$ iff $x$ is a stable semantic fixed point.
\end{axiom}

\begin{axiom}[Middle-way minimization]
$\mathrm{MW}(P)$ minimizes $E$ over $P$ and satisfies the
non-oscillation condition.
\end{axiom}

Together these axioms define what may be called \emph{middle-way
logic}: a semantics whose notion of equivalence is defined by middle-way
selection rather than syntactic identity.

% ============================================================
% 5.2 Middle-Way Equivalence
% ============================================================

\subsection{Middle-Way Equivalence}

The cornerstone of EMT-based logic is the middle-way equivalence
operator $\mweq$.

\begin{definition}[Middle-way equivalence]
For expressions $A$ and $B$, we define
\[
   A \mweq B
\]
to mean that $A$ and $B$ evaluate to the same middle-way value under
EMT evaluation.
\end{definition}

This yields a new form of equality that is:

\begin{itemize}
  \item context-sensitive,
  \item dynamic and energy-based,
  \item stable under perturbation,
  \item free from classical diagonal paradoxes.
\end{itemize}

\begin{proposition}[Absence of contradiction]
No pair $(A,B)$ exists such that both $A \mweq B$ and
$A \not\mweq B$ hold under the same emptiness state.
\end{proposition}

Middle-way equivalence therefore establishes a consistent alternative
to classical equality in self-referential or dynamically instantiated
systems.

% ============================================================
% 6. Conclusion and Future Directions
% ============================================================

\section{Conclusion and Future Directions}

This paper presented EMT (Emergent Middle Test) as a computational and
logical framework based on four interacting layers:
emptiness, ranki, dependent origination, and the middle way.
The resulting system offers:

\begin{itemize}
  \item a mechanism for generating non-deterministic proposals,
  \item an energy-based structure for evaluating semantic coherence,
  \item an operator for selecting dynamically stable fixed points,
  \item a new equivalence relation $\mweq$ that resolves classical
        paradoxes and non-middle fallacies.
\end{itemize}

EMT serves as an alternative foundational framework for logic,
computation, and the formal evaluation of intelligent behavior. Its
interaction with classical computability theory suggests that some
limitations---such as those imposed by diagonal arguments and fixed
point instability---may be reinterpreted rather than accepted as
inevitable.

\subsection{Prospects for applications}

Several promising directions emerge:

\begin{enumerate}
  \item \textbf{Computability theory:}
        EMT provides new fixed-point principles that avoid diagonal
        paradoxes, suggesting new viewpoints on traditionally
        undecidable problems.
  \item \textbf{Complexity theory:}
        The middle-way equivalence $\mweq$ has already been applied to
        reinterpret the $P$ vs.\ $NP$ problem, replacing static
        equivalence with dynamic convergence.
  \item \textbf{Number theory:}
        Similar reinterpretations apply to conjectures such as the
        ABC conjecture and potentially aspects of the Riemann
        hypothesis.
  \item \textbf{Machine learning:}
        EMT generalizes energy-based models with dynamic domain
        generation and oscillation-based regularizers.
  \item \textbf{Foundations of logic:}
        EMT proposes a new equality relation capable of absorbing
        self-reference without contradiction.
\end{enumerate}

\subsection{Outstanding questions}

Several foundational questions remain open:

\begin{itemize}
  \item What are the minimal conditions for $E$ to guarantee the
        existence of a stable middle-way value?
  \item Can middle-way equivalence be embedded into a categorical
        framework?
  \item Does EMT admit a complete proof theory?
  \item Are there systems where classical equality is strictly more
        expressive than middle-way equivalence?
\end{itemize}

Answering these questions will help determine the ultimate scope and
limitations of EMT as a universal framework.

% ============================================================
% Acknowledgments
% ============================================================

\section*{Acknowledgments}

The author thanks the broader research community exploring the
connections between computation, logic, energy-based models, and
foundational principles inspired by cross-cultural philosophy.
Discussions involving alternative formulations of equality and the
role of dynamic semantics in self-referential systems have
substantially shaped this work.

% ============================================================
% References (skeleton form)
% ============================================================

\begin{thebibliography}{99}

\bibitem{Turing1936}
A.~M.~Turing.
\newblock On computable numbers, with an application to the Entscheidungsproblem.
\newblock \emph{Proceedings of the London Mathematical Society}, 1936.

\bibitem{Godel1931}
K.~G\"odel.
\newblock \"Uber formal unentscheidbare S\"atze der Principia Mathematica und verwandter Systeme I.
\newblock \emph{Monatshefte f\"ur Mathematik und Physik}, 1931.

\bibitem{Minsky1967}
M.~Minsky.
\newblock \emph{Computation: Finite and Infinite Machines}.
\newblock Prentice-Hall, 1967.

\bibitem{Hopfield1982}
J.~Hopfield.
\newblock Neural networks and physical systems with emergent collective computational abilities.
\newblock \emph{PNAS}, 1982.

\bibitem{LeCun2006}
Y.~LeCun, S.~Chopra, and R.~Hadsell.
\newblock A tutorial on energy-based learning.
\newblock 2006.

\bibitem{Chiba2025Godel}
M.~Chiba.
\newblock Non-middle fallacy and the re-formulation of G\"odel's incompleteness theorem.
\newblock arXiv preprint, 2025.

\bibitem{Chiba2025PvsNP}
M.~Chiba.
\newblock A middle-way equivalence approach to the $P$ vs.\ $NP$ problem.
\newblock arXiv preprint, 2025.

\bibitem{Chiba2025ABC}
M.~Chiba.
\newblock Re-formulating the ABC conjecture under middle-way equivalence.
\newblock arXiv preprint, 2025.

\end{thebibliography}

\end{document}
